\section{Experiments}
\label{sec:experiments}

\paragraph{Question.} We test whether local orbit consistency reduces symmetry-induced out-of-distribution error and equivariance defect for an unchanged neural-operator backbone. The main experiment uses a periodic \(64\times64\) 2D Navier--Stokes vorticity operator with Galilean boosts. The base model is the same FNO2d architecture in every row; orbit consistency changes training only and does not add an inference-time module.

\paragraph{Protocol.} We compare four methods: supervised FNO, supervised Lie augmentation, supervised orbit consistency, and augmentation plus orbit consistency. The headline low-label setting uses 2\% labeled data and five seeds. The compute-matched augmentation baseline uses six update steps per epoch, while augmentation plus orbit consistency uses four update steps per epoch with \(\lambda_{\mathrm{orb}}=0.10\). All reported rows use the same test-time model architecture and are regenerated by \texttt{python scripts/reproduce\_paper\_artifacts.py}.

\begin{table*}[t]
\centering
\small
\caption{N64 Galilean headline at 2\% labels over five seeds. All rows use the same FNO2d architecture with 2,668,609 parameters. Lower is better for relative \(L^2\), orbit OOD relative \(L^2\), and equivariance defect.}
\label{tab:n64-headline}
\resizebox{\textwidth}{!}{%
\begin{tabular}{lrrrrr}
\toprule
Method & Steps/epoch & Relative \(L^2\) & Orbit OOD \(L^2\) & Eq. defect & Latency ms/sample \\
\midrule
FNO & 4 & \(0.9287\pm0.0228\) & \(0.9568\pm0.0223\) & \(0.6988\pm0.0173\) & \(16.53\pm2.13\) \\
FNO + orbit & 4 & \(0.4646\pm0.0394\) & \(0.4952\pm0.0348\) & \(0.3359\pm0.0170\) & \(16.11\pm1.66\) \\
FNO + aug. & 6 & \(0.5009\pm0.0590\) & \(0.5276\pm0.0525\) & \(0.3816\pm0.0322\) & \(18.76\pm7.34\) \\
FNO + aug. + orbit & 4 & \(\mathbf{0.3676\pm0.0360}\) & \(\mathbf{0.3950\pm0.0334}\) & \(\mathbf{0.2677\pm0.0174}\) & \(16.87\pm3.19\) \\
\bottomrule
\end{tabular}
}
\end{table*}

\begin{table*}[t]
\centering
\small
\caption{Training compute accounting for the headline 2\% N64 Galilean runs. Forward/backward counts and transformed-sample counts are computed from the logged optimizer steps; wall-clock minutes are estimated from immutable RNG checkpoint timestamps for these legacy cached runs, and new runs now write \texttt{train\_wall\_seconds} directly. Relative train time is normalized to the compute-matched supervised augmentation baseline.}
\label{tab:n64-compute}
\resizebox{\textwidth}{!}{%
\begin{tabular}{lrrrrrrrr}
\toprule
Method & Opt. steps & Model fwds. & Bwd. passes & Aug. labeled samples & Orbit samples & Train min. & Rel. train time & Latency ms/sample \\
\midrule
FNO & 600 & 600 & 600 & 0 & 0 & \(14.6\pm0.4\) & \(0.71\times\) & \(16.53\pm2.13\) \\
FNO + orbit & 600 & 1200 & 600 & 0 & 9600 & \(18.7\pm0.6\) & \(0.91\times\) & \(16.11\pm1.66\) \\
FNO + aug. & 900 & 1800 & 900 & 14400 & 0 & \(20.6\pm2.5\) & \(1.00\times\) & \(18.76\pm7.34\) \\
FNO + aug. + orbit & 600 & 1800 & 600 & 9600 & 9600 & \(21.0\pm2.6\) & \(1.02\times\) & \(16.87\pm3.19\) \\
\bottomrule
\end{tabular}
}
\end{table*}

\begin{table}[t]
\centering
\small
\caption{Label-efficiency comparison for compute-matched augmentation versus augmentation plus orbit consistency with \(\lambda_{\mathrm{orb}}=0.10\).}
\label{tab:n64-label-efficiency}
\begin{tabular}{llrrr}
\toprule
Labels & Method & Orbit OOD \(L^2\) & Eq. defect & Seeds \\
\midrule
1\% & Aug. & \(0.7421\pm0.0838\) & \(0.5080\pm0.0456\) & 3 \\
1\% & Aug. + orbit & \(\mathbf{0.6406\pm0.0508}\) & \(\mathbf{0.3999\pm0.0197}\) & 3 \\
2\% & Aug. & \(0.5276\pm0.0525\) & \(0.3816\pm0.0322\) & 5 \\
2\% & Aug. + orbit & \(\mathbf{0.3950\pm0.0334}\) & \(\mathbf{0.2677\pm0.0174}\) & 5 \\
5\% & Aug. & \(0.3237\pm0.0013\) & \(0.2427\pm0.0014\) & 3 \\
5\% & Aug. + orbit & \(\mathbf{0.2515\pm0.0112}\) & \(\mathbf{0.1621\pm0.0078}\) & 3 \\
\bottomrule
\end{tabular}
\end{table}

\begin{table*}[t]
\centering
\small
\caption{Privileged relative-boost calibration at 2\% labels. Direct OOD evaluates the model on the boosted input. The privileged calibration uses the sampled relative boost from the evaluation protocol and the paired in-distribution prediction, so it is not an observable-frame analytic baseline. Lower is better.}
\label{tab:n64-privileged-calibration}
\resizebox{\textwidth}{!}{%
\begin{tabular}{rrrrr}
\toprule
Max boost & Aug. direct OOD & Aug. priv. calib. & Aug. + orbit direct OOD & Aug. + orbit priv. calib. \\
\midrule
0.10 & \(0.5050\pm0.0574\) & \(0.5009\pm0.0590\) & \(\mathbf{0.3715\pm0.0349}\) & \(0.3676\pm0.0360\) \\
0.20 & \(0.5176\pm0.0543\) & \(0.5009\pm0.0590\) & \(\mathbf{0.3840\pm0.0331}\) & \(0.3676\pm0.0360\) \\
0.25 & \(0.5277\pm0.0523\) & \(0.5009\pm0.0590\) & \(\mathbf{0.3943\pm0.0318}\) & \(0.3676\pm0.0360\) \\
0.35 & \(0.5555\pm0.0476\) & \(0.5009\pm0.0590\) & \(\mathbf{0.4242\pm0.0285}\) & \(0.3676\pm0.0360\) \\
0.50 & \(0.6158\pm0.0384\) & \(0.5009\pm0.0590\) & \(\mathbf{0.4910\pm0.0216}\) & \(0.3676\pm0.0360\) \\
\bottomrule
\end{tabular}
}
\end{table*}

\begin{table}[t]
\centering
\small
\caption{Correlation between equivariance defect and orbit OOD relative \(L^2\) across cached N64 Galilean seeds, label fractions, orbit weights, and boost severities. Partial Pearson \(r\) residualizes source, method, label fraction, maximum boost, and \(\lambda_{\mathrm{orb}}\).}
\label{tab:n64-defect-correlation}
\begin{tabular}{lrrrr}
\toprule
Scope & Obs. & Pearson \(r\) & Spearman \(\rho\) & Partial \(r\) \\
\midrule
All & 90 & 0.851 & 0.838 & 0.779 \\
Train/lambda/label & 40 & 0.990 & 0.990 & 0.976 \\
OOD severity & 50 & 0.761 & 0.754 & 0.516 \\
\bottomrule
\end{tabular}
\end{table}

\begin{table}[t]
\centering
\small
\caption{Orbit-weight ablation at 2\% labels. The final setting \(\lambda_{\mathrm{orb}}=0.10\) gives the best completed N64 Galilean result.}
\label{tab:n64-lambda}
\begin{tabular}{rrrr}
\toprule
\(\lambda_{\mathrm{orb}}\) & Orbit OOD \(L^2\) & Eq. defect & Seeds \\
\midrule
0.01 & \(0.5768\pm0.0378\) & \(0.3955\pm0.0235\) & 3 \\
0.05 & \(0.4518\pm0.0397\) & \(0.3161\pm0.0247\) & 5 \\
0.10 & \(\mathbf{0.3950\pm0.0334}\) & \(\mathbf{0.2677\pm0.0174}\) & 5 \\
\bottomrule
\end{tabular}
\end{table}

\begin{table}[t]
\centering
\small
\caption{Relative reductions of augmentation plus orbit consistency over compute-matched augmentation across Galilean OOD severities at 2\% labels.}
\label{tab:n64-severity-reductions}
\begin{tabular}{rrr}
\toprule
Max boost & OOD reduction & Defect reduction \\
\midrule
0.10 & 26.4\% & 33.6\% \\
0.20 & 25.8\% & 31.3\% \\
0.25 & 25.3\% & 30.2\% \\
0.35 & 23.6\% & 27.8\% \\
0.50 & 20.3\% & 23.6\% \\
\bottomrule
\end{tabular}
\end{table}


\begin{figure*}[t]
\centering
\includegraphics[width=0.88\textwidth]{2d_galilean_n64_label_efficiency_lambda_0p1.pdf}
\caption{N64 Galilean label efficiency. Error bars show standard deviation across seeds and faint points show individual seeds. Augmentation plus orbit consistency improves orbit OOD error and equivariance defect at 1\%, 2\%, and 5\% labels.}
\label{fig:n64-label-efficiency}
\end{figure*}

\begin{figure*}[t]
\centering
\includegraphics[width=0.88\textwidth]{2d_galilean_n64_2pct_ood_severity_lambda_0p1.pdf}
\caption{OOD severity at 2\% labels. The dashed line marks the training boost radius. Augmentation plus orbit consistency remains better than compute-matched augmentation from below the training radius through twice the training radius.}
\label{fig:n64-ood-severity}
\end{figure*}

\begin{figure}[t]
\centering
\includegraphics[width=0.96\linewidth]{2d_galilean_n64_defect_ood_correlation.pdf}
\caption{Equivariance defect predicts Galilean orbit OOD error across cached N64 seeds, label fractions, orbit weights, and boost severities. Circles are trained-model evaluations and squares are severity evaluations.}
\label{fig:n64-defect-correlation}
\end{figure}

\paragraph{Main N64 Galilean result.} Table~\ref{tab:n64-headline} shows that at 2\% labels, augmentation plus orbit consistency improves both paper-facing symmetry metrics over compute-matched supervised augmentation. Orbit OOD relative \(L^2\) decreases from \(0.5276\pm0.0525\) to \(0.3950\pm0.0334\), a 25.1\% reduction. Equivariance defect decreases from \(0.3816\pm0.0322\) to \(0.2677\pm0.0174\), a 29.8\% reduction. Relative to the plain supervised FNO baseline, the reductions are 58.7\% for orbit OOD error and 61.7\% for equivariance defect. Parameter counts are identical across the headline rows.

\paragraph{Compute accounting.} Table~\ref{tab:n64-compute} reports training compute for the same headline runs. The compute-matched augmentation baseline uses 900 optimizer steps and 1800 estimated model forward passes, while augmentation plus orbit consistency uses 600 optimizer steps and the same 1800 estimated model forward passes because each step evaluates the original, augmented, and orbit-transformed inputs. Inference latency remains within run-to-run noise because the deployed map is still the same FNO2d backbone.

\paragraph{Label efficiency.} Figure~\ref{fig:n64-label-efficiency} and Table~\ref{tab:n64-label-efficiency} show that the benefit persists across the low-label curve. Against compute-matched augmentation, \(\lambda_{\mathrm{orb}}=0.10\) reduces orbit OOD error by 13.7\%, 25.1\%, and 22.3\% at 1\%, 2\%, and 5\% labels respectively. The corresponding equivariance-defect reductions are 21.3\%, 29.8\%, and 33.2\%. This supports the label-efficiency framing: the orbit objective is most useful when supervision is scarce, while preserving the same inference map.

\paragraph{OOD severity.} Figure~\ref{fig:n64-ood-severity} and Table~\ref{tab:n64-severity-reductions} evaluate Galilean boosts from below the training radius to twice the training radius. The augmentation-plus-orbit model maintains a lower orbit OOD error at every tested boost. The relative OOD-error reductions are 26.4\%, 25.8\%, 25.3\%, 23.6\%, and 20.3\% at maximum boosts 0.10, 0.20, 0.25, 0.35, and 0.50. Equivariance-defect reductions follow the same trend, from 33.6\% at boost 0.10 to 23.6\% at boost 0.50.

\paragraph{Canonicalization calibration.} Table~\ref{tab:n64-privileged-calibration} reports a privileged relative-boost calibration, not a deployable analytic canonicalization baseline. In the OOD severity protocol, each test input \(a\) is paired with a sampled relative boost \(\delta\), producing \(T_\delta a\). The calibration uses this sampled \(\delta\) and the paired in-distribution prediction \(\mathcal G_\theta(a)\), then reports \(T_\delta \mathcal G_\theta(a)\). This is useful for measuring how much of the direct OOD error comes from motion along the tested orbit, but it assumes access to the paired source input and sampled relative transform. It should therefore not be confused with the observable-frame analytic baseline, which reads the absolute ambient velocity \(c\) from the input channels and canonicalizes a single test input. At maximum boost 0.50, augmentation plus orbit consistency has direct OOD error \(0.4910\pm0.0216\), improving substantially over augmentation \(0.6158\pm0.0384\) while remaining above its privileged calibration value \(0.3676\pm0.0360\).

\paragraph{Defect as predictor.} Figure~\ref{fig:n64-defect-correlation} and Table~\ref{tab:n64-defect-correlation} show that equivariance defect is not merely an auxiliary metric. Across 90 cached N64 Galilean observations, defect correlates strongly with orbit OOD relative \(L^2\) (Pearson \(r=0.851\), Spearman \(\rho=0.838\)); after residualizing source, method, label fraction, maximum boost, and \(\lambda_{\mathrm{orb}}\), the partial Pearson correlation remains \(r=0.779\). The relationship also holds within the OOD severity curve alone (\(r=0.761\)), supporting the mechanism that improved orbit consistency tracks improved OOD behavior.

\paragraph{Ablations and scope.} Table~\ref{tab:n64-lambda} reports the orbit-weight ablation at 2\% labels and supports the final \(\lambda_{\mathrm{orb}}=0.10\) setting. Semi-supervised joint orbit training is useful as an appendix result because it reduces equivariance defect, but it is not the strongest predictive-error result. We do not include a 3D PDE in the main study: a same-resolution 3D analogue would multiply raw field size by 64 and introduce a new solver and validation burden. The current contribution is deliberately narrow: exact-symmetry validation plus a reproducible 2D Galilean low-label result.
